\section{Signal Specifications}\label{app:signals}

This appendix provides the full sub-signal architecture for each of the five
detection signals in the A2A fraud-detection framework. All AUC values reported
are from Phase~5 evaluation on the cleaned dataset (665 agents, 1{,}069 human
counterparties).

%----------------------------------------------------------------------
\subsection{Signal 1: Economic Rationality (ER)}
\label{app:sig:er}

Economic Rationality quantifies whether an address's transaction pattern is
consistent with purposeful economic activity or exhibits characteristics of
autonomous agent behavior.

\paragraph{Sub-signals.}
\begin{enumerate}[leftmargin=*]
  \item \textbf{Utility maximisation detection.}
    Identifies circular value flows using depth-limited DFS on the transaction
    graph. A cycle of length $\leq k$ through an address indicates potential
    wash trading or self-dealing.
  \item \textbf{Purpose coherence.}
    Scores each address's economic rationale by measuring the diversity and
    consistency of its counterparty set and transaction purposes.
  \item \textbf{Concentration patterns.}
    Computes a Herfindahl-like index over outgoing value distribution:
    \begin{equation}
      H_i = \sum_{j} \left(\frac{v_{ij}}{\sum_k v_{ik}}\right)^2,
    \end{equation}
    where $v_{ij}$ is the total value sent from address $i$ to counterparty
    $j$. High concentration suggests programmatic routing.
\end{enumerate}

\paragraph{Key observation.}
Observed agent transaction values range from \$0.06 to \$0.91, well below any
reasonable human cognitive cost threshold for initiating a financial
transfer---a strong prior for autonomous execution.

\paragraph{Phase~5 performance.} $\auc = 0.5152$ (cleaned).

%----------------------------------------------------------------------
\subsection{Signal 2: Network Topology (NT)}
\label{app:sig:nt}

Network Topology is the strongest individual signal. It detects structural
anomalies in the transaction graph that are characteristic of agent swarms and
automated enumeration.

\paragraph{Sub-signals and weights.}
\begin{enumerate}[leftmargin=*]
  \item \textbf{Sender centrality} (weight 0.50).
    Measures the out-degree and eigenvector centrality of the address as a
    sender. High-burst out-degree indicates automated enumeration of
    counterparties.
  \item \textbf{Receiver centrality} (weight 0.30).
    Captures in-degree spikes consistent with sink patterns (value
    aggregation points).
  \item \textbf{Path anomaly} (weight 0.20).
    Flags repeated edges (history extraction) and star topologies (agent
    armies controlled by a single orchestrator).
\end{enumerate}

\paragraph{Detected patterns.}
\begin{itemize}[leftmargin=*]
  \item \emph{High out-degree bursts}: rapid enumeration of new counterparties.
  \item \emph{Repeated edges}: the same directed edge observed many times,
    indicating automated history extraction.
  \item \emph{Star topology}: a single hub address connected to many leaf
    addresses with minimal inter-leaf connectivity.
  \item \emph{Sink patterns}: addresses that only receive and never send,
    serving as value collection endpoints.
\end{itemize}

\paragraph{Evolution.}
Phase~3 specification included four sub-components (with Sybil detection and
Louvain community analysis). These were reduced to three in Phase~4
implementation after Sybil/Louvain proved computationally expensive without
proportionate discriminative gain on the ERC-8004 dataset.

\paragraph{Phase~5 performance.} $\auc = 0.5990$ (cleaned)---strongest signal.

%----------------------------------------------------------------------
\subsection{Signal 3: Value Flow (VF)}
\label{app:sig:vf}

Value Flow tracks how value moves through the network, detecting layering and
velocity patterns associated with automated fund movement.

\paragraph{Sub-signals and weights.}
\begin{enumerate}[leftmargin=*]
  \item \textbf{Flow velocity} (weight 0.50).
    Measures the speed at which value transits through an address (time
    between first inflow and subsequent outflow).
  \item \textbf{Layering indicator} (weight 0.50).
    Detects progressive splitting or aggregation of values across
    multiple hops, a hallmark of obfuscation layering.
\end{enumerate}

\paragraph{Removed sub-signal.}
\texttt{net\_flow\_imbalance} was included in Phase~4 but removed after Phase~5
evaluation. It fired paradoxically: 57.1\% of human addresses triggered the
sub-signal versus only 5.6\% of known agents. The imbalance arose because human
counterparties naturally maintain non-zero net flows (they buy more than they
sell, or vice versa), whereas agents exhibited near-zero net flow due to
symmetric automated routing.

\paragraph{Threshold evolution.}
The detection threshold was re-optimised from 0.24 (Phase~4) to 0.08 (Phase~5)
after the sub-signal removal.

\paragraph{Phase~5 performance.} $\auc = 0.5220$ (cleaned).

%----------------------------------------------------------------------
\subsection{Signal 4: Temporal Consistency (TC)}
\label{app:sig:tc}

Temporal Consistency captures timing anomalies in transaction sequences. It
operates bidirectionally, detecting both machine-speed automation and
suspiciously precise human mimicry.

\paragraph{Sub-signals and weights.}
\begin{enumerate}[leftmargin=*]
  \item \textbf{Circadian violation} (weight 0.35).
    Flags addresses active during hours inconsistent with any single
    timezone's waking period.
  \item \textbf{Timing regularity} (weight 0.35).
    Measures the coefficient of variation (CV) of inter-transaction
    intervals. Both extremes are suspicious:
    \begin{itemize}
      \item $\cvt < 0.05$: ``too perfect'' timing mimicry (agents
        attempting to appear human);
      \item $\cvt > 3.0$: erratic bursting characteristic of
        event-driven bots.
    \end{itemize}
  \item \textbf{Burst detection} (weight 0.30).
    Identifies clusters of transactions with inter-arrival times $< 400$\,ms,
    below human reaction thresholds.
\end{enumerate}

\paragraph{Empirical observations.}
Real agents on the ERC-8004 dataset exhibit $\cvt \approx 1.87$, while the most
sophisticated mimicry agents show $\cvt < 0.005$---four orders of magnitude
more regular than genuine agent traffic.

\paragraph{TC-canary property.}
TC degrades fastest among all signals as adversaries adapt. Its Phase~5
$\auc = 0.4647$ (cleaned, below random baseline) confirms that temporal
features are the first to be gamed, consistent with the theoretical prediction
that timing is the cheapest behavioral dimension to manipulate.

%----------------------------------------------------------------------
\subsection{Signal 5: Cross-Platform Correlation (CP)}
\label{app:sig:cp}

Cross-Platform Correlation is architecturally complete but currently inactive
(assigned weight 0.0) because the evaluation dataset contains only Base chain
transactions.

\paragraph{Design.}
CP detects simultaneous presence of the same logical agent across multiple
EVM-compatible chains by matching:
\begin{itemize}[leftmargin=*]
  \item \texttt{CREATE2} deployment addresses across Base, Ethereum mainnet,
    and BNB Chain;
  \item Correlated transaction timing across chains (shared clock source);
  \item Identical bytecode fingerprints at matched addresses.
\end{itemize}

\paragraph{Activation requirements.}
CP requires multi-chain transaction data. Once cross-chain A2A commerce data
becomes available, the signal can be activated by setting its weight to a
non-zero value and supplying multi-chain transaction feeds to the pipeline.

%----------------------------------------------------------------------
\subsection{Signal Weight Evolution}

\begin{table}[h]
\centering
\caption{Signal weight evolution across development phases.}
\label{tab:weight-evolution}
\begin{tabular}{lccc}
\toprule
\textbf{Signal} & \textbf{Phase 3--4} & \textbf{Phase 5} & \textbf{Phase 5 AUC} \\
\midrule
Economic Rationality (ER) & 0.20 & AUC-prop. & 0.5152 \\
Network Topology (NT)     & 0.20 & AUC-prop. & 0.5990 \\
Value Flow (VF)            & 0.20 & AUC-prop. & 0.5220 \\
Temporal Consistency (TC)  & 0.20 & AUC-prop. & 0.4647 \\
Cross-Platform (CP)        & 0.20 & 0.00      & ---    \\
\bottomrule
\end{tabular}
\end{table}

After Phase~5 evaluation, signal weights were updated from uniform
($w_i = 0.20$) to AUC-proportional for the four active signals, with CP
held at zero pending multi-chain data availability.
